% Unidad U006D. Registro dodecafásico integral y descomposición 1+5+6.

\chapter{Registre dodécaphasique intégral, extension cyclique et décomposition \(1+5+6\)}
\label{chap:registro-dodecafasico-integral}

La période observable de cent huit positions contient douze secteurs de
neuf positions. Cette égalité de cardinalités ne suffit pas à décrire
l’état temporel : il faut distinguer le support ordonné, la loi cyclique,
l’indice sectoriel et le registre enrichi des transports, orientations,
retenues et incidences. Cette unité construit les quatre objets et démontre
la reconstruction entière des douze coordonnées de mémoire.

\section{Support ordonné et coordonnée cyclique}

Pour \(m\in\{1,\ldots,12\}\), nous définissons
\begin{equation}
 E_m:=\{9(m-1)+1,\ldots,9m\}.
 \label{eq:u006d-bloques-nonadicos}
\end{equation}
Le support ordonné et le groupe cyclique sont
\begin{equation}
 \Gamma_{108}:=
 E_1\sqcup\cdots\sqcup E_{12},
 \qquad
 \Theta_{108}:=\mathbb Z/108\mathbb Z.
 \label{eq:u006d-gamma-theta}
\end{equation}
\(\Gamma_{108}\) conserve la position, l’ordre et l’appartenance à un bloc ;
\(\Theta_{108}\) conserve la loi de translation. Ce ne sont pas le même objet.

Pour \(\theta\in\Theta_{108}\), soit
\(\widetilde\theta\in\{0,\ldots,107\}\) son représentant. Nous définissons
\begin{equation}
 \beta(\theta):=
 \left[
 \left\lfloor\frac{\widetilde\theta}{9}\right\rfloor
 \right]_{12},
 \qquad
 r(\theta):=1+(\widetilde\theta\bmod9).
 \label{eq:u006d-coordenadas-bloque-residuo}
\end{equation}
La première coordonnée localise le secteur dodécaphasique ; la seconde publie
la position positive \(1,\ldots,9\) à l’intérieur du secteur.

\section{Extension cyclique non scindée}

Soient \(C_n=\mathbb Z/n\mathbb Z\). Les applications
\begin{equation}
 \iota:C_{12}\longrightarrow C_{108},
 \quad \iota([b]_{12})=[9b]_{108},
 \qquad
 p:C_{108}\longrightarrow C_9,
 \quad p([\theta]_{108})=[\theta]_9
 \label{eq:u006d-mapas-extension}
\end{equation}
forment la suite
\begin{equation}
 0\longrightarrow C_{12}
 \xrightarrow{\;\iota\;}C_{108}
 \xrightarrow{\;p\;}C_9
 \longrightarrow0.
 \label{eq:u006d-extension}
\end{equation}
Pour la section ensembliste
\[
 s([r]_9):=[\widetilde r]_{108},
 \qquad 0\leq\widetilde r<9,
\]
le défaut d’additivité est
\begin{equation}
 c(r,r'):=
 \left[
 \left\lfloor
 \frac{\widetilde r+\widetilde r'}9
 \right\rfloor
 \right]_{12},
 \qquad
 s(r)+s(r')
 =s(r+r')+\iota(c(r,r')).
 \label{eq:u006d-cociclo-extension}
\end{equation}

\begin{theorem}[Exactitude, non-scindement et classe cohomologique]
\label{thm:u006d-extension-no-escindida}
La suite \eqref{eq:u006d-extension} est exacte et ne se scinde pas. Pour
l’action triviale,
\begin{equation}
 H^2(C_9,C_{12})
 \simeq C_{12}/9C_{12}
 \simeq C_3,
 \label{eq:u006d-cohomologia-extension}
\end{equation}
et le cocycle \eqref{eq:u006d-cociclo-extension} représente sa classe
génératrice.
\end{theorem}

\begin{proof}
L’image de \(\iota\) est le sous-groupe des multiples de neuf, qui coïncide
avec \(\ker p\) ; \(p\) est surjective. Cela prouve l’exactitude. Si
l’extension se scindait, \(C_{108}\) serait isomorphe à
\(C_{12}\times C_9\), mais
\[
 \exp(C_{108})=108,
 \qquad
 \exp(C_{12}\times C_9)=\operatorname{mcm}(12,9)=36.
\]
L’identité de cocycle est la division euclidienne de
\(\widetilde r+\widetilde r'\) par neuf. Pour un groupe cyclique à action
triviale,
\(H^2(C_9,C_{12})\simeq C_{12}/9C_{12}\) ; la retenue unitaire se projette sur
\([1]\), qui engendre le quotient d’ordre trois.
\end{proof}

Le non-scindement type le retour :
\[
 g_0(\theta+9)=g_0(\theta),
 \qquad
 \beta(\theta+9)=\beta(\theta)+1\pmod{12}.
\]
Neuf itérations ferment la phase visible et translatent d’un secteur ; cent
huit ferment les deux coordonnées de l’horloge. L’état enrichi conserve
en outre la mémoire verticale, le cylindre, la frontière et la généalogie.

\section{Registre temporel enrichi}

\begin{definition}[Registre dodécaphasique orienté]
\label{def:u006d-registro-dodecafasico}
Pour une histoire enrichie, soit \(\tau_m\) la sous-route ordonnée sur
\(E_m\), \(\sigma_m\) son orientation, \(\kappa_m\) la retenue qui franchit la
frontière du bloc et \(\Lambda_m\) son livre local d’incidences. Le
registre complet est
\begin{equation}
 \mathcal R_{12}:=
 \bigl(
 (E_m,\tau_m,\sigma_m,\kappa_m,\Lambda_m)
 \bigr)_{m=1}^{12}.
 \label{eq:u006d-registro-r12}
\end{equation}
\end{definition}

On distingue ainsi quatre niveaux :
\begin{equation}
 \boxed{
 \mathcal R_{12}
 \neq C_{12}
 \neq\Theta_{108}
 \neq\Gamma_{108}.}
 \label{eq:u006d-cuatro-objetos}
\end{equation}
Le premier est un registre enrichi ; le deuxième, un indice de secteur ; le
troisième, un groupe de translations ; le quatrième, un support ordonné. Il existe
des projections entre eux, mais pas d’identifications.

Soit \(q_{\rm mem}\in\mathbb Z_{\geq0}\) le nombre non borné de tours
nonadiques. La coordonnée dodécaphasique n’est que son résidu
\[
 \beta=[q_{\rm mem}]_{12}.
\]
Après cent huit pas,
\[
 \beta\longmapsto\beta,
 \qquad
 q_{\rm mem}\longmapsto q_{\rm mem}+12.
\]
Le calendrier revient ; la mémoire verticale n’est pas réinitialisée.

\section{Extracteur signé d’événements}

Soit \((d_t)_{t=1}^{108}\) le registre de chiffres sur
\(\Gamma_{108}\). Nous définissons
\begin{equation}
 \mathcal D_{\rm evt}:=\{2,4,6,8\},
 \qquad
 \Sigma_{12}:=(+,+,+,-,-,-,+,+,+,-,-,-).
 \label{eq:u006d-datos-eventos}
\end{equation}
Pour \(m=0,\ldots,11\), soit
\begin{equation}
 e_m:=
 \Sigma_{12}(m+1)\,
 \#\{t\in E_{m+1}:d_t\in\mathcal D_{\rm evt}\}
 \in\mathbb Z.
 \label{eq:u006d-eventos-firmados}
\end{equation}
L’application conserve le recensement local et l’orientation sectorielle comme
facteurs séparés.

Nous apparions les deux demi-couronnes par
\begin{equation}
 p_i:=e_i+e_{i+6},
 \qquad
 a_i:=e_i-e_{i+6},
 \qquad 0\leq i\leq5.
 \label{eq:u006d-pares-antipodales}
\end{equation}
Nous définissons le recensement total, les cinq différences par rapport à la sixième paire et
les six antisymétries :
\begin{equation}
 s_C:=\sum_{i=0}^{5}p_i,
 \qquad
 M_i:=p_i-p_5\quad(0\leq i\leq4),
 \qquad
 A_6:=(a_0,\ldots,a_5).
 \label{eq:u006d-coordenadas-uno-cinco-seis}
\end{equation}

\begin{definition}[Lecteur entier \(1+5+6\)]
\label{def:u006d-lector-integral}
Le lecteur linéaire est
\begin{equation}
 \mathcal L_{12}:\mathbb Z^{12}\longrightarrow\mathbb Z^{12},
 \qquad
 \mathcal L_{12}(e)
 :=
 (s_C,M_0,\ldots,M_4,a_0,\ldots,a_5).
 \label{eq:u006d-lector-l12}
\end{equation}
La distribution des coordonnées est
\begin{equation}
 \underbrace{1}_{s_C}
 +\underbrace{5}_{M_0,\ldots,M_4}
 +\underbrace{6}_{a_0,\ldots,a_5}
 =12.
 \label{eq:u006d-uno-cinco-seis}
\end{equation}
\end{definition}

\begin{theorem}[Reconstruction exacte et image entière]
\label{thm:u006d-inversa-l12}
Le lecteur \(\mathcal L_{12}\) est injectif. Sur son image,
\begin{align}
 p_5&=
 \frac{s_C-\sum_{i=0}^{4}M_i}{6},
 &
 p_i&=p_5+M_i\quad(0\leq i\leq4),
 \label{eq:u006d-inversa-p}\\
 e_i&=\frac{p_i+a_i}{2},
 &
 e_{i+6}&=\frac{p_i-a_i}{2}
 \quad(0\leq i\leq5).
 \label{eq:u006d-inversa-e}
\end{align}
Un tuple de sortie appartient à \(\operatorname{im}\mathcal L_{12}\) si et
seulement si
\begin{equation}
 s_C-\sum_{i=0}^{4}M_i\equiv0\pmod6
 \label{eq:u006d-condicion-divisibilidad}
\end{equation}
et, pour les \(p_i\) reconstruits,
\begin{equation}
 p_i\equiv a_i\pmod2
 \qquad(0\leq i\leq5).
 \label{eq:u006d-condicion-paridad}
\end{equation}
\end{theorem}

\begin{proof}
De \(M_i=p_i-p_5\), on obtient
\[
 s_C=\sum_{i=0}^{4}(p_5+M_i)+p_5
 =6p_5+\sum_{i=0}^{4}M_i,
\]
ce qui donne la première équation de \eqref{eq:u006d-inversa-p} ; les autres sont
immédiates. Additionner et soustraire
\(p_i=e_i+e_{i+6}\) et \(a_i=e_i-e_{i+6}\) produit
\eqref{eq:u006d-inversa-e}. La divisibilité par six et les six
congruences de parité sont précisément les conditions nécessaires et
suffisantes pour que toutes les coordonnées reconstruites soient entières.
\end{proof}

La division par six apparaît après le recensement orienté, non pendant celui-ci. Les
conditions d’intégralité de cette unité sont exactement
\eqref{eq:u006d-condicion-divisibilidad} et
\eqref{eq:u006d-condicion-paridad} ; on n’attribue pas ici de forme normale de
Smith à une matrice qui n’aurait pas été préalablement définie.

\section{Deux mots dodécaphasiques de provenances distinctes}

La distinction doit être conservée
\begin{equation}
 U_{12}^{\rm cat}:=U_6\Vert U_6,
 \qquad
 U_{12}^{\rm sgn}:=\mathcal H_{12}(K).
 \label{eq:u006d-dos-palabras-doce}
\end{equation}
La première concatène deux émissions du catalogue de longueur six. La seconde
est une coordonnée signée reliée réversiblement à \(K\) par l’opérateur
\(\mathcal H_{12}\). Leur longueur commune n’implique pas une même provenance,
sémantique ou loi de transport.

\begin{criterion}[Critères de réfutation]
\label{crit:u006d-refutacion}
La construction est réfutée si la suite
\eqref{eq:u006d-extension} n’est pas exacte, si sa classe cohomologique est
triviale, si un tuple de l’image viole
\eqref{eq:u006d-condicion-divisibilidad} ou
\eqref{eq:u006d-condicion-paridad}, si les formules inverses ne retrouvent pas
les douze entiers initiaux, ou si le retour de \(\beta\) est interprété comme
une réinitialisation de \(q_{\rm mem}\).
\end{criterion}
